\chapter{Den tomme formelen}

\begin{motto}
Eit århundre med designutdanning vart bygd kring ein påstand\\ som ikkje kan vere usann, og som difor heller ikkje kan vere sann.
\end{motto}

Kvar disiplin har ein setning ho lærer studentane sine fyrst, og som heile resten heng på. I fysikken er det at kraft er masse gonger akselerasjon. I biologien at form og funksjon heng saman gjennom seleksjon. I designfaget er det at \emph{form følgjer funksjon}. Setninga er over hundre år gamal, ho er tilskriven Sullivan, ho vart Bauhaus\kref{16} sitt slagord og modernismens kjerne, og ho er framleis det næraste faget kjem ein lov. Ho er òg, ved nærare ettersyn, ikkje ein lov i det heile. Ho er ein tautologi forkledd som ein lov, og oppdaginga av at ho er tom er den einaste store sjølvkritiske innsikta faget har produsert om sitt eige fundament. At innsikta er korrekt, og at faget likevel ikkje bygde noko nytt på henne, fortel deg alt du treng å vite om disiplinen sin evne til å lære av seg sjølv.

\section{Kvifor formelen ikkje kan vere usann}

Jan Michl viste det med eit enkelt grep\kref{1}. Spør kva «funksjon» tyder i formelen. To svar er moglege, og begge gjer formelen tom.

Det fyrste svaret er at funksjon tyder \emph{det objektet skal gjere} i snever forstand — stolen skal bere ein sitjande kropp. Men dette underbestemmer forma fullstendig. Tusenvis av radikalt ulike stolar ber den same kroppen like godt. Pinnestolen, lenestolen, krakken, barokkstolen, Eames-stolen, plaststolen i hagen; alle løyser sittefunksjonen, og likevel deler dei knapt ei linje. Om form følgde funksjon i denne tydinga, ville alle stolar konvergert mot éi form. Dei gjer det openbert ikkje. Altså følgjer forma ikkje funksjonen i snever forstand.

Det andre svaret reddar formelen ved å utvide «funksjon» til å tyde alt: ikkje berre sittinga, men det kulturelle uttrykket, statusmarkeringa, den estetiske gleda, produksjonsøkonomien, tidsånda. No stemmer formelen alltid, fordi kva som helst ein form gjer kan kallast funksjonen han følgjer. Barokkstolen følgjer sin funksjon (å syne makt), plaststolen følgjer sin (å vere billeg), designarstolen følgjer sin (å signalisere smak). Men ein påstand som er sann same kva du observerer, fortel deg ingenting før du observerer. Han kan ikkje brukast til å føreseie ei einaste form, kan ikkje skilje ein god form frå ein dårleg, kan ikkje rettleie eit einaste val. Han er ikkje gal; han er tom. Han har forma til ei forklaring og innhaldet til ei sirkelslutning.

Mellom desse to lesingane finst det ingen tredje som både er ikkje-triviell og sann. Anten betyr funksjon noko presist, og då er formelen usann fordi forma er underbestemt; eller funksjon betyr alt, og då er formelen sann men tom. Dette er ikkje eit retorisk poeng. Det er den logiske strukturen til faget sin grunnsetning, og strukturen er den til ein tautologi. Eit heilt århundre med formgjevarar vart oppdregne på ein regel som ikkje kunne brytast, og forveksla det med ein regel som ikkje kunne feile.

\section{Det formelen skjulte}

Det verste er ikkje at formelen var tom. Det verste er kva tomleiken skjulte. Når du seier at forma følgjer funksjonen, og funksjon betyr alt, har du i praksis sagt at forma følgjer \emph{noko} — og latt vere usagt kva, kven som bestemmer det, og kva som vert lagt til side. Formelen gav modernismen eit alibi for å presentere reine estetiske og ideologiske val som naudsynte konsekvensar av funksjonen. Det skråstilte taket, det kvite huset, fråvêret av ornament; alt vart framstilt som det funksjonen kravde, når det i røynda var det ein bestemt smak føretrekte. Michl viser at funksjonalismen ikkje var funksjonell; han var ein stil som hevda å vere fråvêret av stil\kref{19}. Og fordi grunnsetninga gjorde stilen usynleg — kalla han funksjon — kunne faget aldri diskutere han som det han var: eit val, mellom andre moglege val, gjort på vegne av nokon, med konsekvensar for andre.

Her ser du strukturen som går att gjennom heile boka. Faget sin grunnsetning gjer ikkje arbeidet ein grunnsetning skal gjere — å gjere vala synlege og diskuterbare. Han gjer det motsette. Han gøymer vala bak ein retorikk om naudsyn. «Forma følgjer funksjonen» tyder i praksis «mine val er ikkje val, dei er det saka krev». Det er ikkje eit fundament. Det er eit slør.

\section{Hundre år utan oppfølging}

Michl skreiv kritikken sin alt på nittitalet, og bygde på røter tilbake til åttitalet\kref{19}. Han er ikkje obskur; han er sitert, antologisert, undervist. Og her er det avslørande: faget tok imot kritikken, nikka, la han på pensum ved sida av Sullivan-sitatet, og endra ingenting. Det heldt fram med å lære studentane formelen som om kritikken ikkje fanst, og lærte dei kritikken som om han ikkje hadde konsekvensar. Dette er ikkje korleis ein disiplin med eit fundament oppfører seg. Då Michelson og Morley viste at eteren ikkje fanst, bygde fysikken ein ny teori. Då formelen vart vist tom, bygde designfaget ingenting. Det la berre endå eit lag på den voksande bunken av ting faget «veit» men ikkje brukar.

Grunnen er enkel og dømmande. For å erstatte ein tom formel treng du ein ikkje-tom. Du treng å kunne seie, presist, kva forma faktisk følgjer — kva som verkar på henne, korleis, og kvar dei verkande kreftene kjem frå. Det krev eit omgrep om forma som ein posisjon i eit rom av moglegheiter, og om kreftene som gradientar i det rommet, og om formgjevaren som ein agent som navigerer det. Med andre ord krev det presis det traktaten \textsc{Formlære} legg fram, og som ingen i faget hadde lagt fram då Michl skreiv. Faget kunne diagnostisere tomleiken, men hadde ingenting å setje i staden, og difor heldt det fram med å bruke det det visste var tomt. Ein lege som veit at medisinen er verkningslaus, men held fram med å skrive han ut fordi han ikkje har nokon annan, er ikkje lenger ein lege. Han er ein som deler ut flasker.

\vignett

Formelen var den intellektuelle kjernen i designutdanninga, og kjernen var hol. Men ein kunne tenkje at den verkelege kunnskapen i faget aldri budde i formlar uansett; at han budde i praksis, i atelieret, i overføringa frå meister til student. Det er det neste forsvaret, og det er sterkare. La oss gå inn i atelieret.
